\section{Coupled invariants by descending linear spaces}
\label{sec:spaces}

\subsection{Conservation is stronger than induction needs}

An equation $p(x)=0$ can be invariant even if $p(F(x))\ne p(x)$ as a polynomial.
For example, $p(F(x))=2p(x)$ preserves zero. More generally, a family of equations
can preserve its common zero set while its individual components mix:
\begin{equation}\label{eq:matrix-action}
 \mathbf p(F_a(x))=M_a\mathbf p(x),\qquad
 \mathbf p=(p_1,\ldots,p_r)^{\mathsf T},\quad M_a\in\Q^{r\times r}.
\end{equation}
This is a genuine step beyond searching only for conserved quantities.

Forge explicitly recognizes the still more general relation
$p_i(F(x))=\sum_j h_{ij}(x)p_j(x)$ and correctly warns that jointly discovering
unknown $p_j$ and unknown multipliers is bilinear \cite{forge-structural}.
The algorithm in this section does not pretend otherwise. It restricts the
multipliers to constants and computes a greatest stable subspace by linear
preimages, never solving for a new invariant basis and its action matrix in one
bilinear system. Section~\ref{sec:ideals} removes the constant-multiplier
restriction by a different, target-seeded algorithm.

\subsection{Problem statement}

Fix finitely many polynomial maps
$F_1,\ldots,F_s:\Q^n\to\Q^n$ and a polynomial initialization
$\phi:\Q^k\to\Q^n$. Every finite word in the maps is allowed; thus the problem
includes arbitrary switching between transitions. The desired target is an
equation $q(x)=0$ along every resulting orbit.

Let $P_d=\Q[x_1,\ldots,x_n]_{\le d}$. Its dimension is
\[
 N=\binom{n+d}{d}.
\]
The initialization induces a linear map
$E_\phi:P_d\to\Q[t_1,\ldots,t_k]$, $p\mapsto p\circ\phi$.
Each transition induces a linear pullback map $F_a^*:p\mapsto p\circ F_a$.
Pullback is linear in the polynomial being transformed, even when $F_a$ itself
is nonlinear. Its image need not lie in $P_d$.

Define
\begin{align}
 V_0&=\kerop E_\phi,\label{eq:V0}\\
 V_{j+1}&=V_j\cap\bigcap_{a=1}^{s}(F_a^*)^{-1}(V_j).
 \label{eq:descending}
\end{align}
All spaces on the right are subspaces of $P_d$. Membership of a pullback in
$V_j$ is tested as an equality of full polynomials, not after projection back
to degree $d$.

\begin{theorem}[Greatest stable bounded-degree space]\label{thm:space}
Iteration \eqref{eq:descending} stabilizes after at most $N$ strict dimension
decreases. Its fixed point $V_*$ is the greatest subspace of $P_d$ whose
polynomials vanish on $\phi(\Q^k)$ and which is stable under every $F_a^*$.
Every polynomial in $V_*$ vanishes on every reachable state.
\end{theorem}
\begin{proof}
The sequence is descending in the finite-dimensional space $P_d$. A proper
inclusion decreases dimension by at least one; once equality occurs, applying
the same operator again leaves the space unchanged.

At a fixed point, $V_*\subseteq V_0$ and $F_a^*(V_*)\subseteq V_*$ for every
$a$. Let $W\subseteq P_d$ have those two properties. Then $W\subseteq V_0$.
If $W\subseteq V_j$, stability gives $F_a^*W\subseteq W\subseteq V_j$, so
$W\subseteq V_{j+1}$. Hence $W\subseteq V_*$.

Finally, all elements of $V_*$ vanish initially. If they all vanish at $x$,
then for $p\in V_*$, the polynomial $p\circ F_a$ also belongs to $V_*$, so
$p(F_a(x))=0$. Induction on the length of an arbitrary transition word proves
the reachability assertion.
\end{proof}

\subsection{An executable preimage computation}

Choose an ordered monomial basis for $P_d$ and keep $V_j$ in reduced row-echelon
form, with row basis $b_1,\ldots,b_r$. For each $a$ and $i$, expand
$b_i\circ F_a$ exactly. Eliminate its coefficients at the pivot monomials of
$V_j$ using the basis rows, retaining all remaining monomials. Denote the
remainder by $u_{a,i}$.

A combination $p=\sum_i c_i b_i$ has its pullback in $V_j$ exactly when
\[
 \sum_i c_i u_{a,i}=0
 \quad\text{for each }a.
\]
Stack the coefficient equations for all remainder monomials and transitions.
Their nullspace gives the allowed vectors $c$. Multiply that nullspace basis by
the old row basis and row-reduce once more. This computes
\eqref{eq:descending} with ordinary exact linear algebra.

\begin{lstlisting}
V := nullspace(coefficients(monomialBasis o initialization))
repeat:
  residuals := []
  for transition F:
    for basis polynomial b in V:
      residuals.append(fullRemainder(b o F, vectorSpace = V))
  C := simultaneousCoefficientNullspace(residuals)
  next := canonicalRowSpace(C * V)
  if next = V: break
  V := next
return V, exactActionMatrices(V), targetCoordinates(V)
\end{lstlisting}

The phrase \emph{full remainder} is load-bearing. If $x'=x^2$ and $d=1$,
then $x\circ F=x^2$ is not a degree-one polynomial. Replacing it by zero would
falsely make $\Span\{x\}$ look stable. The prototype retains the $x^2$
coefficient and rejects this proposed stability. Failure to find a stable
space is a search limitation, not a proof that $x=0$ fails on the orbit.

\subsection{A coupled nonlinear example}

Take the parameterized initial state
\[
 (x,y,z)=(t,t^2,2t^2),\qquad t\in\Q,
\]
and allow either transition
\begin{align*}
 F_1(x,y,z)&=(2x,\;2y+z,\;y+3z+x^2),\\
 F_2(x,y,z)&=(-x,\;2x^2-y,\;y+z-x^2).
\end{align*}
We want to prove $z=2y$ after every finite sequence of these updates.

Set $e_1=y-x^2$ and $e_2=z-2x^2$. Direct expansion gives
\[
 \begin{pmatrix}e_1\circ F_1\\e_2\circ F_1\end{pmatrix}
 =\begin{pmatrix}2&1\\1&3\end{pmatrix}
 \begin{pmatrix}e_1\\e_2\end{pmatrix},\qquad
 \begin{pmatrix}e_1\circ F_2\\e_2\circ F_2\end{pmatrix}
 =\begin{pmatrix}-1&0\\1&1\end{pmatrix}
 \begin{pmatrix}e_1\\e_2\end{pmatrix}.
\]
The target is $e_2-2e_1$. Neither equation needs to remain numerically unchanged
outside its zero set; only the pair's common zero set needs to be preserved.

The degree-two search has 10 monomials. Initialization leaves a five-dimensional
vanishing space, and one strict refinement reduces it to the two-dimensional
stable space generated by $e_1,e_2$. The stored canonical basis happens to use
the order $(e_2,e_1)$; its action matrices are correspondingly permuted. The
conservation-only ablation, with the same degree and initialization, has dimension
zero. A degree-one search also has zero stable dimension and misses the target.

These are exact nullspace calculations, not fits to sampled trajectories.
Initialization is checked after substituting the symbolic parameter $t$, and
each step relation is checked coefficient by coefficient.

\subsection{A stronger completeness statement for affine transitions}

For nonlinear transitions, $V_*$ may omit valid degree-$d$ invariants because
preserving them can require higher-degree multiples. If every $F_a$ is affine,
then $P_d$ is itself pullback-stable, and a sharper statement holds.

\begin{theorem}[Affine-fragment completeness]\label{thm:affine}
For affine transitions and polynomial initialization over $\Q$, $V_*$ is exactly
the set of degree-at-most-$d$ polynomials that vanish on all reachable states.
\end{theorem}
\begin{proof}
Let $W$ be that set. It is a subspace of $P_d$ and vanishes initially.
If $p\in W$, then $p\circ F_a$ has degree at most $d$. For any reachable $x$,
$F_a(x)$ is also reachable because every transition is enabled, so
$(p\circ F_a)(x)=0$. Thus $W$ is pullback-stable and is contained in $V_*$ by
Theorem~\ref{thm:space}. The reverse inclusion is its soundness conclusion.
\end{proof}

The prototype tests this theorem by a second, forward formulation. Let
$m(x)$ be the column of all monomials through degree $d$. For an affine map,
$m(F_a(x))=L_a m(x)$ for an exact rational matrix $L_a$. Start from the span of
coefficient vectors of the parameterized feature column $m(\phi(t))$, and close
that span under all $L_a$. Its annihilator is the space of polynomial relations.

Why use coefficient vectors rather than only a few initial feature values?
A rational linear functional vanishes on all evaluations of a polynomial vector
if and only if the resulting polynomial is identically zero, which over the
infinite field $\Q$ is equivalent to every coefficient vanishing. Thus the two
spans have the same annihilator. Forward reachable-span closure and backward
stable-space refinement must therefore agree.

They agree on all 72 generated affine test systems. Of these, 24 have a nonzero
relation space; the others are useful zero-space controls, not additional
nontrivial theorems. The two algorithms share the exact scalar-arithmetic helper,
so this is algorithmic differential testing, not two wholly independent
implementations of rational arithmetic.

\subsection{The certificate and its proof cost}

A positive certificate contains the polynomials $p_i$, one constant matrix
$M_a$ per transition, and target coordinates $q=\sum_i c_i p_i$. The checker
verifies
\begin{equation}\label{eq:space-certificate}
 p_i\circ\phi=0,\qquad
 p_i\circ F_a=\sum_j (M_a)_{ij}p_j,\qquad
 q=\sum_i c_i p_i.
\end{equation}
It neither recomputes the fixed point nor trusts the searcher's claim of
maximality. Linear independence of the listed polynomials is also unnecessary
for certificate soundness; redundant generators merely cost space.

A Lean implementation can express the invariant as a finite conjunction of
polynomial equations and use a single generic reachability-induction lemma.
For small certificates, generating direct expressions and discharging the
identities with \code{ring} or \code{ring\_nf} is the simplest bridge. For larger
ones, a reflected sparse-polynomial checker should use an established polynomial
semantics such as \code{MvPolynomial.eval\textsubscript{2}}, whose definition and evaluation
lemmas are present at the pinned Mathlib revision \cite{mathlib-eval}.

The expensive work is not the theorem above but coefficient growth and
substitution. At degree $d$, a transition of degree $D$ can produce monomials
through degree $dD$, with ambient count $\binom{n+dD}{n}$. There are at most
$N$ strict refinements, but this bound does not make the arithmetic polynomial
in the input bit length or keep the intermediate polynomials small. Cache
monomial pullbacks, exploit sparse rows, and use fraction-free or modular exact
linear algebra in production. Python-FLINT supplies rational row reduction and
integer nullspace facilities that can serve the proposer; acceptance still
uses the exact identities \cite{flint-q,flint-z}.

\subsection{Guards and control locations}

Ignoring guards enlarges the transition relation. An invariant proved for the
enlarged system is valid for the guarded system, but failure in the enlarged
system says nothing negative about the guarded source program. Parameters in
$\phi$ are likewise universally unconstrained in this fragment.

A useful unimplemented extension is one space $V_\ell$ per control location.
For an edge $e:\ell\to m$, require $F_e^*V_m\subseteq V_\ell$, with separate
initial constraints at entry locations. Simultaneous descending refinement in
$\prod_\ell P_d$ still terminates after at most $|L|N$ strict dimension drops.
The certificate becomes a family of rectangular action matrices, one per edge.
This is the natural representation for mutually recursive helpers whose
invariants differ by phase; it is preferable to artificially embedding every
control location in the same polynomial equation list.
